\documentclass[11pt,a4paper]{article}
\usepackage[utf8]{inputenc}
\usepackage[margin=1in]{geometry}
\usepackage{amsmath,amssymb}
\usepackage{booktabs}
\usepackage{graphicx}
\usepackage{hyperref}
\usepackage{cite}
\usepackage{listings}
\usepackage{color}

\hypersetup{
    colorlinks=true,
    linkcolor=blue,
    filecolor=magenta,      
    urlcolor=cyan,
    pdftitle={Personalized Movie Recommendation System},
    pdfpagemode=FullScreen,
}

\title{\textbf{FlixPulse: A Personalized Movie Discovery \& Recommendation System}\\
\large{Cult Open Projects 2026 --- AI/ML Section}}
\author{\textbf{IIT Roorkee Cultural Council}}
\date{\today}

\begin{document}

\maketitle

\begin{abstract}
Personalized content discovery is crucial for user retention and engagement on modern digital entertainment platforms. This technical report details the implementation of FlixPulse, a movie recommendation system trained on a dense, high-quality subset of the Netflix Prize Dataset containing 5.7 million ratings. We implement, compare, and serialize two recommendation models: Singular Value Decomposition (SVD) and Item-Based K-Nearest Neighbors (KNN) Collaborative Filtering. The models are evaluated using prediction-based metrics (RMSE, MAE) and ranking-oriented metrics (MAP@10, Precision@10, Recall@10, NDCG@10, Hit Rate, and Intra-List Diversity). Our results demonstrate that SVD achieves superior prediction accuracy (RMSE = 0.8160) and ranking performance (NDCG@10 = 0.8337) with high efficiency. Furthermore, we develop an item-similarity explanation framework and cold-start fallbacks, deploying the complete pipeline as a full-stack FastAPI and React web application.
\end{abstract}

\newpage

\section{Problem Understanding \& Motivation}
Recommendation systems are the cornerstone of digital content providers like Netflix, YouTube, and Spotify. By learning user preferences from historical interaction logs, these systems deliver personalized recommendations that directly influence user retention, search times, and platform engagement.

The core technical challenges addressed in this system include:
\begin{itemize}
    \item \textbf{Data Sparsity}: In real-world datasets, users interact with a tiny fraction of the available items, leading to highly sparse interaction matrices ($>99\%$ sparsity).
    \item \textbf{Prediction vs. Ranking Trade-off}: While regression models minimize prediction error (RMSE), user satisfaction relies on ranking accuracy (getting relevant items at the top of the feed).
    \item \textbf{Explainability}: Providing clear rationale for recommendations (e.g., \textit{``Because you gave Toy Story a 5 stars''}) increases transparency and user trust.
    \item \textbf{Cold Start}: Handling recommendations when new users register with no ratings, or when new movies are uploaded with zero interaction records.
\end{itemize}

---

\section{Exploratory Data Analysis (EDA)}
We analyzed and prepared a dense subset of the Netflix Prize Dataset to construct our modeling workflows. 

\subsection{Sparsity and Dataset Dimensions}
The final filtered subset has the following dimensions:
\begin{itemize}
    \item \textbf{Total Ratings ($|R|$)}: 5,702,695 reviews
    \item \textbf{Unique Users ($|U|$)}: 21,045 active users (minimum 200 ratings)
    \item \textbf{Unique Movies ($|M|$)}: 1,287 popular movies (minimum 2,000 ratings)
    \item \textbf{Density}: 21.04\% (Matrix Sparsity: 78.96\%)
\end{itemize}

\subsection{Rating Frequency Distribution}
Ratings are mapped on a 1--5 star scale. The frequency distribution demonstrates a significant positive skew:
\begin{table}[h]
\centering
\caption{Rating Distribution Summary}
\begin{tabular}{ccc}
\toprule
\textbf{Rating Value} & \textbf{Count} & \textbf{Percentage (\%)} \\
\midrule
1★ & 183,184 & 3.21\% \\
2★ & 432,610 & 7.59\% \\
3★ & 1,489,612 & 26.12\% \\
4★ & 2,096,180 & 36.76\% \\
5★ & 1,501,109 & 26.32\% \\
\bottomrule
\end{tabular}
\label{tab:rating_dist}
\end{table}

The mean rating of the dataset is \textbf{3.73}, which indicates a positive selection bias. Users are naturally more inclined to watch and rate movies they expect to enjoy.

\subsection{User Activity \& Movie Popularity}
The distribution of ratings per user shows a heavy-tailed profile, where the median user has rated 234 movies, but power users reach up to 1,500 ratings. Similarly, blockbuster titles capture a large fraction of total reviews (e.g., \textit{Miss Congeniality} and \textit{Independence Day} having over 10,000 reviews in our subset).

---

\section{Methodology \& Model Design}
We implemented and validated two distinct recommendation paradigms to resolve the user-item prediction problem.

\subsection{Singular Value Decomposition (SVD) Matrix Factorization}
SVD maps both users and items to a joint latent factor space of dimensionality $f$. The user-item rating prediction is formulated as:
\begin{equation}
\hat{r}_{u,i} = \mu + b_u + b_i + p_u^T q_i
\end{equation}
Where:
- $\mu$ represents the global mean rating.\\
- $b_u$ and $b_i$ represent user and item biases respectively.\\
- $p_u \in \mathbb{R}^f$ is the latent factor vector for user $u$.\\
- $q_i \in \mathbb{R}^f$ is the latent factor vector for item $i$.

The parameters are optimized using Stochastic Gradient Descent (SGD) with L2 regularization to minimize the squared error:
\begin{equation}
\min_{b, p, q} \sum_{(u,i) \in K} (r_{u,i} - \hat{r}_{u,i})^2 + \lambda (b_u^2 + b_i^2 + \|p_u\|^2 + \|q_i\|^2)
\end{equation}
Where $\lambda$ is the regularization term to prevent overfitting.

\subsection{Item-Based Collaborative Filtering (KNN with Means)}
A neighborhood-based algorithm that calculates similarities between movies. The prediction is calculated as:
\begin{equation}
\hat{r}_{u,i} = \mu_u + \frac{\sum_{j \in N_u(i)} \text{sim}(i, j) \cdot (r_{u,j} - \mu_j)}{\sum_{j \in N_u(i)} |\text{sim}(i, j)|}
\end{equation}
Where:
- $\mu_u$ is the user $u$'s mean rating.\\
- $N_u(i)$ is the set of items rated by user $u$ that are most similar to item $i$.\\
- $\text{sim}(i, j)$ is the cosine similarity matrix weight.

---

\section{Evaluation Metrics \& Results}
The models were trained on 80\% of the data and tested on the remaining 20\%. We computed both mandatory metrics and several optional validation criteria.

\subsection{Metric Formulations}
\begin{itemize}
    \item \textbf{RMSE} (Root Mean Squared Error): Measures prediction error magnitude.
    \item \textbf{MAE} (Mean Absolute Error): Average absolute difference.
    \item \textbf{MAP@10} (Mean Average Precision @ 10): Measures rank accuracy of relevant items (actual rating $\ge 3.5$).
    \item \textbf{NDCG@10} (Normalized Discounted Cumulative Gain): Graded position relevance weight.
    \item \textbf{Diversity (Intra-List Diversity)}: Average difference between items:
    \begin{equation}
    ILD@10 = \frac{2}{k(k-1)}\sum_{i \in R} \sum_{j \in R, j \neq i} (1 - \text{sim}(i, j))
    \end{equation}
\end{itemize}

\subsection{Experimental Results}
\begin{table}[h]
\centering
\caption{Model Performance Evaluation Summary}
\begin{tabular}{lcc}
\toprule
\textbf{Metric} & \textbf{SVD Model} & \textbf{Item-Based KNN} \\
\midrule
\textbf{Test RMSE} (Lower is better) & \textbf{0.8160} & 0.8553 \\
\textbf{Test MAE} (Lower is better) & \textbf{0.6368} & 0.6687 \\
\textbf{MAP@10} (Higher is better) & \textbf{0.7427} & 0.7079 \\
\textbf{Precision@10} (Higher is better) & \textbf{0.8051} & 0.7814 \\
\textbf{Recall@10} (Higher is better) & \textbf{0.3388} & 0.3272 \\
\textbf{NDCG@10} (Higher is better) & \textbf{0.8337} & 0.8081 \\
\textbf{Hit Rate@10} & \textbf{99.94\%} & 99.92\% \\
\textbf{User Coverage} & \textbf{99.96\%} & 99.96\% \\
\textbf{Intra-List Diversity} & \textbf{0.0477} & 0.0452 \\
\textbf{Training Time (s)} & \textbf{8.92s} & 225.86s \\
\textbf{Inference Time (s)} & \textbf{0.31s} & 4.85s \\
\bottomrule
\end{tabular}
\label{tab:results}
\end{table}

---

\section{Recommendation Examples \& Explainability}
SVD generates highly accurate recommendations but acts as a black box. To provide explainability, we couple SVD predictions with the KNN similarity weights.

For user \texttt{134001}, the system generates:
\begin{enumerate}
    \item \textbf{Movie}: \textit{Toy Story 2} (Predicted: 4.8★)
    \item \textbf{Explanation}: \textit{``Because you gave 'Toy Story' a high rating of 5 (Similarity: 98\%)''}
\end{enumerate}

---

\section{Key Insights \& Performance Trade-offs}
\begin{itemize}
    \item \textbf{Prediction vs. Ranking Correlation}: SVD's superior rating prediction accuracy (RMSE 0.8160) directly drives better rank placement (NDCG@10 = 83.37\%), outperforming neighbor-based filtering across all vectors.
    \item \textbf{Complexity vs. Latency}: KNN requires calculating large similarity weight tables, increasing training times to 225.8s. SVD represents latent states linearly, updating in 8.9s.
\end{itemize}

---

\section{Future Improvements}
\begin{itemize}
    \item \textbf{Deep Learning Integrations}: Transitioning to Neural Collaborative Filtering (NCF) to model non-linear interactions.
    \item \textbf{Scraping External Metadata}: Enhancing UI dashboard aesthetics by scraping TMDB API poster and trailer backdrops.
\end{itemize}

\end{document}
